% =====================================================================
% config/preamble.tex  --  GÓI LỆNH VÀ CẤU HÌNH CHUNG
% Trình biên dịch: XeLaTeX hoặc LuaLaTeX (KHÔNG dùng pdfLaTeX).
% =====================================================================
\RequirePackage{iftex}
\ifPDFTeX
  \PackageError{preamble}{Dang dung pdfLaTeX}{Hay doi trinh bien dich sang XeLaTeX hoac LuaLaTeX (Overleaf: Menu > Compiler).}
\fi

\usepackage[a4paper,top=\MarginTop,bottom=\MarginBottom,left=\MarginLeft,right=\MarginRight]{geometry}
\usepackage[fontsize=\BodyPt]{fontsize}

% --- Ngôn ngữ và font (Unicode) ---------------------------------------
\usepackage{fontspec}
\usepackage{anyfontsize} % cho phép cỡ chữ tùy ý (tránh cảnh báo cỡ 13pt của phông toán)
\usepackage[vietnamese]{babel}
\babelfont{rm}{\MainFont}
\babelfont{sf}{TeX Gyre Heros}
% Font đơn cách cho mã nguồn: TeX Gyre Cursor có sẵn trong TeX Live/Overleaf và có đủ dấu tiếng Việt.
% (Một số bản DejaVu Sans Mono thiếu glyph tiếng Việt, nên không dùng làm mặc định.)
\babelfont{tt}{TeX Gyre Cursor}

% Việt hóa tên danh mục cho đồng bộ (babel mặc định: "Danh sách hình vẽ", "Danh sách bảng").
\addto\captionsvietnamese{%
  \renewcommand{\listfigurename}{Danh mục hình}%
  \renewcommand{\listtablename}{Danh mục bảng}}

% --- Gói thông dụng -----------------------------------------------------
\usepackage{amsmath}
\usepackage{xcolor}
\usepackage{graphicx}
\graphicspath{{figures/}{uml/}}
\usepackage{booktabs,longtable,array,tabularx,multirow}
\usepackage{enumitem}
\setlist{itemsep=2pt,topsep=4pt}
\usepackage{setspace}
\usepackage{comment}
\usepackage{float}
\usepackage{caption}
\captionsetup{font=small,labelfont=bf,labelsep=period}
% Quy ước: chú thích bảng đặt TRÊN bảng, chú thích hình đặt DƯỚI hình.

% --- Hộp màu (dùng cho hộp hướng dẫn) -----------------------------------

% --- Mã nguồn ----------------------------------------------------------
\usepackage{listings}
\lstset{
  basicstyle=\ttfamily\footnotesize, breaklines=true, columns=fullflexible,
  keepspaces=true, frame=single, numbers=left, numberstyle=\tiny\color{gray},
  tabsize=2, showstringspaces=false, captionpos=b}
\renewcommand{\lstlistingname}{Mã nguồn}
\renewcommand{\lstlistlistingname}{Danh mục mã nguồn}

% --- Tiêu đề, giãn dòng, đầu/chân trang --------------------------------
\usepackage{titlesec}
\titleformat{\chapter}[hang]{\normalfont\bfseries\Large}{\chaptername\ \thechapter.}{0.5em}{}
\titleformat{name=\chapter,numberless}[hang]{\normalfont\bfseries\Large}{}{0pt}{}
\titlespacing*{\chapter}{0pt}{0pt}{22pt}
\titleformat{\section}{\normalfont\bfseries\large}{\thesection.}{0.5em}{}
\titleformat{\subsection}{\normalfont\bfseries}{\thesubsection.}{0.5em}{}
\titleformat{\subsubsection}{\normalfont\bfseries\itshape}{\thesubsubsection.}{0.5em}{}
\setcounter{secnumdepth}{3}
\setcounter{tocdepth}{2}
\setstretch{\LineStretch}
\setlength{\parindent}{1cm}
\setlength{\parskip}{3pt}

\usepackage{fancyhdr}
\setlength{\headheight}{15pt}
\renewcommand{\chaptermark}[1]{\markboth{\chaptername\ \thechapter.\ #1}{}}
\pagestyle{fancy}
\fancyhf{}
\fancyhead[L]{\small\itshape\leftmark}
\fancyfoot[C]{\thepage}
\renewcommand{\headrulewidth}{0.4pt}
\fancypagestyle{plain}{\fancyhf{}\fancyfoot[C]{\thepage}\renewcommand{\headrulewidth}{0pt}}

% --- Liên kết (hyperref phải đứng trước cleveref; cleveref đứng cuối) -----
\usepackage[unicode]{hyperref}
\hypersetup{colorlinks=true,linkcolor=blue!45!black,citecolor=blue!45!black,urlcolor=blue!45!black}
% In ra giấy: đổi dòng trên thành  \hypersetup{hidelinks}
\usepackage[nameinlink]{cleveref}
\crefname{figure}{Hình}{Hình}\Crefname{figure}{Hình}{Hình}
\crefname{table}{Bảng}{Bảng}\Crefname{table}{Bảng}{Bảng}
\crefname{chapter}{Chương}{Chương}\Crefname{chapter}{Chương}{Chương}
\crefname{section}{Mục}{Mục}\Crefname{section}{Mục}{Mục}
\crefname{subsection}{Mục}{Mục}\Crefname{subsection}{Mục}{Mục}
\crefname{subsubsection}{Mục}{Mục}\Crefname{subsubsection}{Mục}{Mục}
\crefname{appendix}{Phụ lục}{Phụ lục}\Crefname{appendix}{Phụ lục}{Phụ lục}
\crefname{equation}{Công thức}{Công thức}\Crefname{equation}{Công thức}{Công thức}
\crefname{lstlisting}{Mã nguồn}{Mã nguồn}\Crefname{lstlisting}{Mã nguồn}{Mã nguồn}
